% !TEX root = ../main.tex
\clearpage
\phantomsection
\section*{5 \textbf{연구 방법론}}
\label{sec:research-methodology}
\addcontentsline{toc}{section}{5 연구 방법론}

본 절은 EndorseRank 구현, Adaptive Weighted PageRank(AWP)와의 비교, 평판 구조·거래 성공 정렬·계산 효율성 측면에서의 성능 평가를 위한 전반적 접근을 개괄한다.


\phantomsection
\subsection*{5.1 연구 설계 및 방법}
\label{sec:research-design-and-method}
\addcontentsline{toc}{subsection}{5.1 연구 설계 및 방법}

\textbf{데이터셋 개요}


본 연구는 공개 온체인 이벤트 로그를 주요 데이터셋으로 사용한다. 원시 데이터는 Arbitrum One 스마트 계약이 방출한 타임스탬프, 주소 수준 기록으로 재현 가능하나 본질적으로 가명이다(지갑은 식별자이지 검증된 실세계 주체가 아님).


\textbf{핵심 분석 단위}


지갑 주소 및 소유자에서 지출자로의 방향(일방향) 관계, 이후 희소 그래프로 표현된다. GMX V2 거래 결과는 \textbf{PositionDecrease} 이벤트의 \textbf{account} 필드를 사용하여 동일 지갑 수준에서 집계된다.


\textbf{주요 신호}


소유자, 지출자, 승인 금액을 지정하는 ERC-20 Approve(및 EIP-2612 Permit) 승인(각 승인의 맥락은 토큰 계약이 제공). GMX V2 \textbf{PositionIncrease} 및 \textbf{PositionDecrease} 이벤트는 검증을 위한 실현 거래 결과를 제공한다.


\textbf{희소성 및 heavy tail}


보증/승인 그래프는 차수/가중치 분포가 치우친(sparse and skewed) 고도로 희소할 것으로 예상되며(소수 계약/라우터가 많은 승인을 받음), 평균 기반 요약을 넘어선 견고한 평가를 동기 부여한다.


\textbf{시간 구조}


이벤트는 시계열을 형성한다; EndorseRank는 고정 관찰 기간(2025년 12월 1일--2026년 5월 31일)에서 운영되며, 관계당 가장 최근 관련 승인 정보를 보존하여 ``최신 상태'' 엣지 가중치를 구성한다. GMX 거래 성공 대리변수는 동일 기간에서 계산된다.


\textbf{데이터 품질 고려}


이벤트 로그는 일반적으로 신뢰할 수 있으나 취소(금액 = 0), ``기본 전액 지출'' 승인, 간헐적 인덱싱 공백을 포함할 수 있으며, 이러한 특성은 전처리 및 견고성 검사에서 명시적으로 다룬다. GMX 이벤트는 EventEmitter \textbf{EventLog1} 형식에서 디코딩해야 하며, 트레이더 신원은 라우터 또는 keeper 주소 오귀속을 피하기 위해 \texttt{transaction.from}이 아닌 이벤트 \textbf{account} 필드를 사용한다.


데이터 수집 및 전처리


\textbf{블록체인 데이터 소스}


\textbf{Arbitrum One}(체인 ID 42161), 관찰 기간 2025년 12월 1일--2026년 5월 31일, Google BigQuery(\texttt{bigquery-public-data.goog\_blockchain\_arbitrum\_one\_us}) 또는 아카이브 노드 경유.


이벤트 추출


승인/permit 파생 승인: Arbitrum One의 모든 ERC-20 Approval(owner, spender, value) 이벤트, 토큰이 표준 Approval 이벤트를 방출하는 permit 호환 흐름을 통해 생성된 승인 포함.


GMX V2 무기한 스왑 결과: GMX V2 \textbf{EventEmitter} 계약(\texttt{0xC8ee91A54287DB53897056e12D9819156D3822Fb})의 \textbf{PositionDecrease} 로그, \textbf{account}, \textbf{basePnlUsd}, \textbf{orderType}, 청산 플래그 추출을 위해 디코딩. 수익성 청산은 \textbf{basePnlUsd} $> 0$이며 포지션 청산이 아님(\texttt{orderType} $\neq 7$).


\textbf{실증 검증 선택}


이론적으로 온체인 신용 평판에 대한 가장 강한 검증 대상은 담보 부족 대출(예: health-factor 위반 후 대출 청산)에서 추출한 이진 부도 또는 청산 레이블이다. 그러한 레이블은 온체인에서 신용 실패를 직접 인코딩한다. 그러나 실무적으로 본 연구 대상 Arbitrum One 생태계는 관찰 기간 내 대규모 재현 가능한 지갑 수준 부도 데이터셋을 구성하기에 충분한 활성 공개 무담보 또는 담보 부족 대출 거래량이 부족하다. 다른 체인에서 비교 가능 레이블을 수집하면 단일 체인 설계를 깨고 보증 그래프 의미론 개선 없이 데이터 엔지니어링 비용을 증가시킨다.


따라서 본 학위논문은 \textbf{GMX V2 무기한 롱/숏 포지션 결과}를 주요 실증 검증 틀로 채택한다. GMX의 담보화된 무기한 거래는 실현 PnL이 있는 고빈도 지갑 귀속 \textbf{PositionDecrease} 이벤트를 생성하여, 보증 그래프와 동일 체인에서 신용 유사 리스크 관리에 대한 측정 가능한 대리변수---청산 성공 횟수, 실현 이익 대리변수, 청산 성공률---를 가능하게 한다. 이 측정치는 \textbf{도메인 직관 대리변수}로 해석되며 ground-truth 대출 상환 레이블이 아니다; 한계는 6절에서 논의한다.


\textbf{거래 성공 대리변수(검증 대상)}


Do et al.(2023)이 PageRank를 in-degree 및 in-value에 대해 검증한 것을 따르면, GMX 유사치를 다음과 같이 정의한다:

\begin{itemize}
\item[\textbullet] \textbf{청산 성공 횟수:} 관찰 기간 내 지갑당 수익성 \textbf{PositionDecrease} 이벤트 수.
\item[\textbullet] \textbf{실현 이익 대리변수:} 기간 내 수익성 청산에 대한 $\sum \max(\text{basePnlUsd}, 0)$.
\item[\textbullet] \textbf{청산 성공률:} 승/전체 청산, 통계적 안정성을 위해 $\text{total\_closes} \geq 3$인 지갑만 유지.
\end{itemize}


표본 추출


정렬 분석은 평판 그래프에 존재하고 관찰 기간에 GMX \textbf{PositionDecrease} 이벤트가 최소 3회 있는 지갑의 교집합을 사용한다. 계산 벤치마크는 고정 시드로 100,000 지갑의 재현 가능한 무작위 표본을 사용하며, 견고성 검사는 50,000 및 200,000 지갑 표본을 비교하여 확장성 민감도를 검증한다.


\begin{longtable}{>{\raggedright\arraybackslash}p{0.18\linewidth}>{\raggedright\arraybackslash}p{0.39\linewidth}>{\raggedright\arraybackslash}p{0.35\linewidth}}
\toprule
\textbf{구성개념} & \textbf{정의} & \textbf{관련 변수} \\
\midrule
평판\newline 구조 & 순위 알고리즘이 산출하는 지갑 평판의 패턴 및 상대적 순서. & (1) EndorseRank 점수;\newline AWP 점수;\newline EndorseRank와 AWP 간 순위 상관(Spearman's \emph{\(\rho\)} / Kendall's \emph{\(\tau\)}) \\
\addlinespace
거래 성공\newline 정렬 & 평판 점수가 GMX V2 거래 성공 대리변수 순위와 정렬되는 정도. & 청산 성공 횟수;\newline 실현 이익 대리변수;\newline 청산 성공률;\newline 각 순위 방법 대 각 대리변수에 대한 Spearman's \emph{\(\rho\)} / Kendall's \emph{\(\tau\)} \\
\addlinespace
계산\newline 효율성 & 점수 산출에 필요한 계산 자원(시간, 메모리, 수렴 반복). & EndorseRank 런타임;\newline AWP 런타임;\newline 최대 메모리 사용량;\newline 수렴까지 반복 \\
\addlinespace
\bottomrule
\end{longtable}

\proposalSubheading{1. 그래프 구성 및 알고리즘 구현}

\begin{itemize}
\item[\textbullet] \textbf{Approve/Permit 그래프 }($G_{\text{approve}}$)
\begin{itemize}
\item[\textendash] \textbf{노드}: 지갑 주소.
\item[\textendash] 엣지: u에서 v로의 최신 0이 아닌 승인으로 가중된 방향 소유자-지출자 관계. ERC-20 토큰은 decimals와 경제적 규모가 다르므로, 가중치는 먼저 토큰 decimal 조정; 민감도 검사는 원시 승인 가중치, log 변환 가중치, 신뢰할 수 있는 가격 데이터가 있는 경우 가치 정규화 가중치를 비교한다.
\end{itemize}
\item[\textbullet] \textbf{송금 기반 그래프 }($G_{\text{transfer}}$)
\begin{itemize}
\item[\textendash] 동일 노드 집합; \emph{u }\(\to\) \emph{v} 송금 금액 합으로 가중된 엣지.
\end{itemize}
\item[\textbullet] 알고리즘
\end{itemize}
\textbf{EndorseRank: }$G_{\text{approve}}$에 대한 \textbf{가중 PageRank, 감쇠 인자 d = 0.85.} 0.85는 표준 PageRank 설정, 확산과 텔레포테이션의 안정적 균형, 선행 그래프 순위 연구와의 직접 비교를 위해 유지된다.


\begin{itemize}
\item[\textendash] \textbf{AWP 재현}: $G_{\text{transfer}}$에 대한 Adaptive Weighted PageRank
\end{itemize}
Do et al.(2023)에 따라 시간 및 가치 감쇠 포함.


\begin{itemize}
\item[\textbullet] 도구 및 환경
\begin{itemize}
\item[\textendash] NetworkX 또는 GraphFrames를 사용한 Python; Scipy 희소 솔버; 재현성을 위한 Docker 컨테이너.
\end{itemize}
\end{itemize}
\proposalSubheading{2. 성능 평가 절차}

\begin{itemize}
\item[\textbullet] 거래 성공 정렬
\begin{itemize}
\item[\textendash] 각 순위 방법(EndorseRank, AWP, 송금 기반 PageRank)에 대해 매칭 지갑 표본에서 청산 성공 횟수, 실현 이익 대리변수, 청산 성공률 순위 대 Spearman's \emph{\(\rho\)} 및 Kendall's \emph{\(\tau\)} 계산.
\item[\textendash] EndorseRank 대 AWP 정렬 계수 비교; EndorseRank가 1.0에 접근하지 않으면서(단순 대리변수 순위와의 중복을 나타냄) 더 높은 상관을 달성하는지 검증.
\item[\textendash] Do et al.(2023)과 유사한 표로 결과 보고, 단순(비가중) 송금 PageRank 기준선과의 비교 포함.
\end{itemize}
\item[\textbullet] 순위 발산
\begin{itemize}
\item[\textendash] EndorseRank와 AWP 간 Spearman's \emph{\(\rho\)} 및 Kendall's \emph{\(\tau\)} 계산.
\item[\textendash] 산점도 및 top-k Jaccard 겹침으로 시각화.
\end{itemize}
\item[\textbullet] 계산 효율성
\begin{itemize}
\item[\textendash] 하드웨어: 고정 22-vCPU, 64 GB RAM 서버.
\item[\textendash] 측정:
\end{itemize}
\end{itemize}
\(\ast\) 런타임(그래프 구축 + PageRank 수렴, 5회 실행 평균)


\(\ast\) 최대 메모리 사용량


\(\ast\) 허용오차 10\textsuperscript{-}\textsuperscript{6} 도달까지 반복


\phantomsection
\subsection*{5.2 데이터 분석}
\label{sec:data-analysis}
\addcontentsline{toc}{subsection}{5.2 데이터 분석}

\proposalSubheading{1. 가설 검증}

\begin{itemize}
\item[\textbullet] \textbf{H1 (Divergence)}
\begin{itemize}
\item[\textendash] EndorseRank와 AWP 간 순위 상관이 1과 유의미하게 다른지 일표본 t-검정 또는 Wilcoxon 부호순위 검정으로 검증.
\end{itemize}
\item[\textbullet] H2 (거래 성공 정렬)
\begin{itemize}
\item[\textendash] 각 GMX 대리변수에 대해 EndorseRank 대 AWP의 Spearman's \emph{\(\rho\)} 및 Kendall's \emph{\(\tau\)} 계수를 쌍별 부트스트랩 신뢰구간 또는 상관된 상관에 대한 Fisher z 변환 검정으로 비교.
\end{itemize}
\item[\textbullet] H3 (계산 효율성)
\begin{itemize}
\item[\textendash] 런타임, 메모리, 반복 횟수에 대한 쌍별 t-검정(또는 비정규 시 Wilcoxon 검정).
\end{itemize}
\end{itemize}
\proposalSubheading{2. 견고성 검사}

\begin{itemize}
\item[\textbullet] 감쇠 인자 \emph{d }\(\in\) \{0\emph{.}75\emph{, }0\emph{.}85\emph{, }0\emph{.}95\} 변화.
\item[\textbullet] Arbitrum One 거래량 상위 ERC-20 토큰 부분그래프에서 분석 반복.
\item[\textbullet] 표본 크기(50k 대 200k 지갑) 민감도 평가.
\item[\textbullet] 원시 대 크기 가중 실현 이익 대리변수 정의 비교.
\end{itemize}
\proposalSubheading{3. 재현성}

\begin{itemize}
\item[\textbullet] 모든 코드, 쿼리, 분석 스크립트를 GitHub에 버전 관리.
\item[\textbullet] 데이터 스냅샷 및 Docker 이미지 보존으로 결과 전체 재현.
\end{itemize}
이 방법론적 틀은 평판 구조, GMX 거래 성공 정렬, 계산 성능 측면에서 EndorseRank가 기존 PageRank 기반 접근보다 갖는 이점에 대한 엄밀하고 재현 가능한 평가를 보장한다.
